\section{\acrfull{HAOT}\label{ch3:haot}}
An open-source tool, \gls{HAOT}, has been developed to perform \gls{AO}
analysis using \gls{CFD} results as input. This tool enables automated
post-processing of flowfield data to extract relevant optical metrics in
hypersonic nonequilibrium regimes, aiding in the analysis and prediction of
aero-optical distortions. \gls{HAOT} is published on
\href{https://pypi.org/project/haot/}{PyPI}, with its source code hosted on
\href{https://github.com/mliza/HAOT}{GitHub}, and its documentation provided on
\href{https://haot.readthedocs.io/en/latest/}{Read the Docs}. The current
version $1.1.2$ of \gls{HAOT} includes five modules:
\begin{itemize} 
\item \textbf{Aerodynamics:} Implements commonly used aerodynamic functions.
\item \textbf{Optics:} Provides equations used in aero-optics and hypersonic
    flow analysis.
\item \textbf{Quantum Mechanics:} Offers commonly used quantum mechanical
    functions relevant to hypersonic flow.
\item \textbf{Constants:} Contains various constants, primarily
    spectroscopy-related used in optics.
\item \textbf{Conversion:} Provides a suite of conversion functions used in
    this module. 
\end{itemize}
The \gls{HAOT} directory structure is given on \Cref{fig:haotStructure}.
\begin{figure}[H]
    \centering
    \begin{lstlisting}[style=tree]
                                HAOT
                                ├── LICENSE
                                ├── README.md
                                ├── CHANGELOG.md 
                                ├── pyproject.toml
                                └── requirements.txt 
                                HAOT/haot
                                ├── __init__.py
                                ├── aerodynamics.py
                                ├── optics.py
                                ├── quantum_mechanics.py 
                                ├── constants.py
                                └── conversions.py 
                                HAOT/unit_test
                                ├── __init__.py
                                ├── test_aerodynamics.py
                                ├── test_optics.py
                                └── test_quantum_mechanics.py
                                    \end{lstlisting}
    \caption{\gls{HAOT} project directory structure}
    \label{fig:haotStructure}
\end{figure}


Validation results are provided in subsequent sections. The constants and
conversion functions were validated through visual inspection. The primary
objective of \gls{HAOT} is to enable researchers to perform optical analysis
using \gls{CFD} results. \gls{HAOT} is an open-source tool designed with the
intention of allowing other developers to expand its current functionalities.
The goal of the tool is to serve as a \textit{software blueprint} that provides
a foundational framework for researchers, particularly in academia, where
software tools are often poorly written, inadequately tested, insufficiently
documented, and underdeveloped \cite{Robillard1998,Liu2008}. Comprehensive
documentation for the tool, including usage instructions with or without
\gls{CFD} data, can be found on the
\href{https://haot.readthedocs.io/en/latest/}{documentation page}. \\

The \gls{HAOT} tool is integrated with a GitHub Actions pipeline that is
automatically triggered whenever a \gls{PR} is submitted. This \gls{CI}
workflow ensures the integrity of the codebase by running automated tests and
checks to verify that new contributions do not introduce errors or regressions.
By validating each proposed change before it is merged, the pipeline plays a
critical role in maintaining the stability, reliability, and reproducibility of
the \gls{HAOT} software. \\

In addition to the unit tests, every time a \gls{PR} is merged into the main
branch, the project’s documentation is automatically regenerated and published
to the official \href{https://haot.readthedocs.io/en/latest/}{documentation
page}. This automated deployment ensures that the documentation remains
synchronized with the latest version of the codebase, providing users and
developers with up-to-date information about the functionality, usage, and
configuration of the \gls{HAOT} tool. The process is managed through a \gls{CD}
pipeline, which integrates with Read the Docs to streamline the publication of
updated content. \\ 

Defensive programming is the core software development technique employed in
the development of the \gls{HAOT} tool. Defensive programming is a software
safety and reliability practice that operates under the assumption that any
input, even from seemingly reliable sources, may be erroneous, malicious, or
unexpected. As such, the program is designed to anticipate and handle incorrect
usage, edge cases, and failures gracefully \cite{Denault2006}. \\

In the context of \gls{HAOT}, this means that every user-provided input—such as
simulation files, physical parameters, numerical data, and file paths—is
rigorously validated before the program proceeds with any operations. Input
validation routines check for correct data types, valid ranges, logical
consistency, and completeness to prevent the software from executing under
invalid or dangerous conditions. This ensures robustness, fault tolerance, and
easier debugging, especially in high-stakes engineering contexts where accuracy
and reliability are critical. \\

\Cref{fig:defensiveProgramming} illustrates the structure and philosophy behind
the function implementation in this work. By design, each function assumes that
the user may provide invalid or nonsensical input—commonly referred to as
"garbage in."\\

The first check is for undefined inputs. For example, if the function prototype
requires a positive integer and the user instead passes a string, the function
will catch this mismatch and output an appropriate error message. \\

The second layer checks for logical errors. Suppose the correct input type is
provided (an integer), but the value is semantically incorrect (e.g., a
negative number instead of a positive one). In that case, the function again
detects the issue and halts execution with an informative error. \\

Only after passing these two layers of input validation—type correctness and
semantic correctness—does the function proceed to the solution space, where the
actual computation or processing is performed.  
\begin{figure}[H]
    \makebox[\textwidth][c]{\includegraphics[width=0.5\textwidth]{chapter3/defensiveProgramming}}
    \caption{Defensive Programmings}\label{fig:defensiveProgramming}
\end{figure}
By adopting a defensive programming paradigm, \gls{HAOT} is more resistant to
crashes, data corruption, and unintended behavior, which is particularly
important in research tools used for high-fidelity scientific computing and
analysis in hypersonic \gls{AO} simulations. Results to the unit test can be found on \crefApendex{app:haotResults}.
